Redistricting is inherently a decadal process, yet most research examines it in isolated snapshots rather than across time. We present the first comprehensive longitudinal analysis of algorithmic redistricting, applying a consistent recursive bisection methodology to census data from 2000, 2010, and 2020. Our analysis reveals three key findings: (1) algorithmic compactness remains stable at 0.45-0.46 Polsby-Popper across two decades despite significant population shifts, (2) enacted district compactness declined by 12\% over the same period, and (3) states adopting redistricting commissions were associated with 7.3 percentage point improvements relative to non-commission states (though causality cannot be definitively established given selection effects). We compute geographic stability metrics using Intersection-over-Union analysis, finding that 61\% of districts maintained substantial boundary continuity from 2010 to 2020. These results establish algorithmic redistricting not as a one-time intervention but as a durable foundation for fair representation across demographic change.
